\section{Numerical checks and the production metric modulus}
\label{sec:continued-numerics}
These computations check specific formulas and finite-grid
behavior. The continuum estimates above follow from the
proved bounds.

\begin{figure}[htbp]
\centering
\begin{tikzpicture}
\begin{axis}[width=.47\textwidth,height=.32\textwidth,
  xlabel={Fraction of the observation window},
  ylabel={Relative orbit displacement},
  ylabel style={font=\scriptsize},xlabel style={font=\small},
  tick label style={font=\scriptsize},
  ymin=-.27,ymax=.28,xmin=0,xmax=1,
  legend style={font=\scriptsize,at={(.5,.03)},anchor=south,
                legend columns=2,draw=none,fill=none}]
\addplot[teal!75!black,thick] table[x=fraction,y=branch_a,col sep=comma]
 {figures/collision_path.csv};
\addlegendentry{First history}
\addplot[orange!85!black,thick,dashed] table[x=fraction,y=branch_b,col sep=comma]
 {figures/collision_path.csv};
\addlegendentry{Second history}
\addplot[gray,densely dotted,forget plot] coordinates {(.2666666667,-.27) (.2666666667,.28)};
\end{axis}
\end{tikzpicture}\hfill
\begin{tikzpicture}
\begin{loglogaxis}[width=.47\textwidth,height=.32\textwidth,
  xlabel={Fourier frequency $k$},ylabel={Normalized orbit gain},
  xlabel style={font=\small},ylabel style={font=\small},
  tick label style={font=\scriptsize},
  legend style={font=\scriptsize,at={(.03,.03)},anchor=south west,draw=none,fill=none}]
\addplot[orange!85!black,thick,mark=square*] table[x=mode,y=nodal,col sep=comma]
 {figures/null_gains.csv};
\addlegendentry{Nodal slice}
\addplot[teal!75!black,thick,mark=*] table[x=mode,y=placed,col sep=comma]
 {figures/null_gains.csv};
\addlegendentry{Off-node slice}
\addplot[gray,densely dotted,forget plot] table[x=mode,y=principal,col sep=comma]
 {figures/null_gains.csv};
\end{loglogaxis}
\end{tikzpicture}
\caption{An exact rank-two collision on the marked torus.
The left panel shows $\widetilde P$ at one angle after
subtracting the common homogeneous background.
The curves match at both ends. The dotted line marks
the replacement terminal slice.
The right panel compares the linearized orbit response
on carriers that are in free resonance at both locations.
The response decays at the original nodal slice and
approaches the nonzero transport prediction (dotted)
at the replacement slice.}
\label{fig:null-numerics}
\end{figure}

For Fig.~\ref{fig:null-numerics}, the coordinate tangent
equations were integrated independently of the null-ray
transfer formula. The two ray angles agreed with
direct line integration within $9\times10^{-16}$.
The off-node normalized gain tends to the predicted
$0.08559$ on the tested sequence, whereas the nodal
gain falls from $0.03133$ at $k=64$ to $0.002041$
at $k=1024$. At bandwidth $96$, two grids gave the
minimum full-orbit singular gain $0.04090$ at the
replacement slice and gains below $5\times10^{-13}$
at the original nodal slice. A separate smaller-amplitude
test verified quadratic variation of the $P$-operator
and linear variation of the $Q$-operator as predicted
in Theorem~\ref{th:exact-rank-recovery}.

For nonlinear inversion at fixed times, an independent
higher-grid Fourier evolution supplied a perturbed
rank-two record. At bandwidth $96$, the two inverse
resolutions agreed to the shown precision.
Adding a normalized record error $3\times10^{-5}$ in
one resonant $P$ mode gave a recovered velocity
$H^1$ error $3.62\times10^{-4}$ from the full orbit
fields. Reading $P$ alone gave an $H^1$ velocity
error $0.06022$ despite a final position residual
below $6\times10^{-11}$. With exact data the full
orbit reconstruction reached a velocity error of
about $5.3\times10^{-10}$.
Reading $P$ and the terminal metric gradient in the
same noisy run gave $5.21\times10^{-6}$ velocity
error in the tested direction. These comparisons use
the stated Fourier-weighted norms and a single, specified
noise direction.

\begin{figure}[htbp]
\centering
\begin{tikzpicture}
\begin{loglogaxis}[width=.68\textwidth,height=.30\textwidth,
  xlabel={Fourier frequency $k$},ylabel={Minimum normalized gain},
  ymin=2e-5,ymax=.29,tick label style={font=\scriptsize},
  xtick={153,467,624,1575,2352},xticklabels={153,467,624,1575,2352},
  legend style={font=\scriptsize,at={(.03,.03)},anchor=south west,draw=none}]
\addplot[orange!85!black,thick,mark=square*] table[x=mode,y=orbit,col sep=comma]
 {figures/full_metric_gains.csv};
\addlegendentry{Orbit fields}
\addplot[teal!75!black,thick,mark=*] table[x=mode,y=metric,col sep=comma]
 {figures/full_metric_gains.csv};
\addlegendentry{Orbit fields and metric gradient}
\end{loglogaxis}
\end{tikzpicture}
\caption{Both velocities are missing in the circular-loop
example at an irrational delay. The plotted modes sample
near cancellations of the orbit-only transfer.
The metric gradient at the same terminal slice
raises the normalized gain.}
\label{fig:irrational-metric}
\end{figure}

Figure~\ref{fig:irrational-metric} compares smallest
singular gains for both velocities at sampled modes on
a circular-loop background at $d=0.02$.
At $k=1575$ the orbit-only gain is about
$4.23\times10^{-5}$ and the gain with the metric
gradient is about $0.118$. At the longer
illustrative gap $d=\pi\sqrt2/24$, the corresponding
gains at $k=3356$ are about $1.4\times10^{-5}$
and $0.678$. These are modal numerical comparisons.
The all-frequency short-delay claim in
Proposition~\ref{prop:irrational-metric} follows
from the analytic reverse estimate.

For the family of Proposition~\ref{prop:from-observations-to-tail},
the computed holonomy angle divided by
$\epsilon^2$ was $-2.94478497$ at
$\epsilon=2.5\times10^{-4}$, against the
quadratic prediction $-2.94478495$.
Two spatial grids gave depth-28 estimates of
the five asymptotic quantities agreeing to
$8\times10^{-14}$ in the stated norms.
The limiting momentum identity had residual
$1.24\times10^{-14}$ in $H^1$ on the finer grid.
These finite-depth checks illustrate the two
steps of the result and the asymptotic constraint.

The frequency-assigned experiment used a separate
rank-two pair for which the two original orbit snapshots
agreed within $1.3\times10^{-12}$ in direct forward
evolution. The delayed sampling resolved their
initial-velocity difference and had a normalized gain
of about $0.891$ on that pair. At bandwidth $288$ four
delayed slots gave a minimum free gain $0.529$,
compared with $2.46\times10^{-4}$ using the baseline
slot alone. The nonlinear reconstruction used an
independent fourth-order spatial stencil. At that
bandwidth it failed to converge on the coarsest
inverse grid and the recovered velocity improved
on the two finer grids. The mismatch is consistent
with stencil dispersion at inadequate spatial
resolution. Unknown timing offsets within the tested
$1/K$ scale introduced visible bias, which was reduced
when the true timestamps were used.
Substituting the exact continuum damped-wave transfer
for the stencil's free transfer removed the principal
dispersion discrepancy. At bandwidth $288$ on the
coarse grid the corrected inverse gave an $H^1$
velocity error $3.60\times10^{-9}$ in two updates
on the small-amplitude test. A more strongly coupled
profile retained a nonzero nonlinear model discrepancy
after the same correction. This separates finite-grid
phase error from nonlinear model error.

\begin{figure}[htbp]
\centering
\begin{tikzpicture}
\begin{semilogyaxis}[width=.47\textwidth,height=.32\textwidth,
  xlabel={Iteration from the secant},ylabel={Normalized record residual},
  xlabel style={font=\small},ylabel style={font=\small},
  tick label style={font=\scriptsize},
  legend style={font=\scriptsize,at={(.97,.98)},anchor=north east,draw=none,fill=none}]
\addplot[teal!75!black,thick,mark=*] table[x=iteration,y=residual,col sep=comma]
 {figures/archive_iteration.csv};
\addlegendentry{Archived grid $N=4096$}
\addplot[gray,densely dashed] table[x=iteration,y=certified_upper,col sep=comma]
 {figures/archive_iteration.csv};
\addlegendentry{Certified contraction bound}
\end{semilogyaxis}
\end{tikzpicture}\hfill
\begin{tikzpicture}
\begin{loglogaxis}[width=.47\textwidth,height=.32\textwidth,
  xlabel={Spatial grid size},ylabel={Absolute error in $c_\lambda$},
  xtick={256,512,1024,2048},xticklabels={256,512,1024,2048},
  xlabel style={font=\small},ylabel style={font=\small},
  tick label style={font=\scriptsize},
  legend style={font=\scriptsize,at={(.98,.54)},anchor=north east,draw=none,fill=none}]
\addplot[orange!85!black,thick,mark=square*] table[x=grid,y=zero_tail_depth8,col sep=comma]
 {figures/modulus_convergence.csv};
\addlegendentry{Earlier anchor, zero tail}
\addplot[teal!75!black,thick,mark=*] table[x=grid,y=bessel_tail_depth8,col sep=comma]
 {figures/modulus_convergence.csv};
\addlegendentry{Earlier anchor, Bessel tail}
\end{loglogaxis}
\end{tikzpicture}
\caption{Left, the archived reconstruction error in the
normalized record during the secant-started iteration,
compared to the factor from Theorem~\ref{th:archive-certificate}.
Right, the depth-eight smooth production metric-modulus
estimate with the finite-depth anchor at $\tau=6$ and
both field snapshots earlier than the discarded endpoint
at depth eight. The error is against the independent
continuum reference.}
\label{fig:archive-modulus}
\end{figure}

On all three archived centered grids the secant-started
replay took nine updates to reach a normalized residual
below $6\times10^{-13}$. At $N=4096$ it returned to the
held-out velocity within $4.1\times10^{-12}$ in
supremum norm. The numerical replay used the same
stencil as the archived tape. The separate interval
certificate guarantees a contraction in a whole velocity
ball for the exact semidiscrete evolution. A second
interval calculation permits simultaneous variation
of the recorded $P,Q$, weighted momentum and
normalized secant, at the stated radii.

The earlier modulus checklist proposed using the
depth-16 versus depth-24 difference as a spatial
convergence test. An independent Fourier--Bessel
calculation from the production initial data gives
$c_\lambda\simeq0.1180324739805121$ in the historical
anchor normalization, with an analytic tail bound
below $2.1\times10^{-28}$ after $\tau=34$.
At depth eight, the pre-endpoint, zero-tail estimate
has error close to $3.79\times10^{-4}$ on the fine grids.
Using the exact Bessel tail with the same earlier
anchor and two numerical field snapshots gives,
over $N=256,512,1024,2048$, errors close to
$8.93\times10^{-6},\,2.23\times10^{-6},\,
5.58\times10^{-7},\,1.40\times10^{-7}$.
The analytic semidiscrete reference reproduces those
spatial errors. The pre-endpoint depth-16 and
depth-24 zero-tail estimates differ by about
$2.3$ to $2.8\times10^{-10}$ over the grids,
without second-order decrease of that difference.
The spatial error common to the two estimates
cancels in the comparison, leaving the finite-depth
bias and time-rounding effects.
It is the errors against the independent continuum
reference that establish spatial convergence.
The tail calculation for this particular production
profile uses its known finite Fourier support;
the metric anchor remains part of the supplied data.
